%Version 3.1 December 2024
% See section 11 of the User Manual for version history
%
%%%%%%%%%%%%%%%%%%%%%%%%%%%%%%%%%%%%%%%%%%%%%%%%%%%%%%%%%%%%%%%%%%%%%%
%%                                                                 %%
%% Please do not use \input{...} to include other tex files.       %%
%% Submit your LaTeX manuscript as one .tex document.              %%
%%                                                                 %%
%% All additional figures and files should be attached             %%
%% separately and not embedded in the \TeX\ document itself.       %%
%%                                                                 %%
%%%%%%%%%%%%%%%%%%%%%%%%%%%%%%%%%%%%%%%%%%%%%%%%%%%%%%%%%%%%%%%%%%%%%

%%\documentclass[referee,sn-basic]{sn-jnl}% referee option is meant for double line spacing

%%=======================================================%%
%% to print line numbers in the margin use lineno option %%
%%=======================================================%%

%%\documentclass[lineno,pdflatex,sn-basic]{sn-jnl}% Basic Springer Nature Reference Style/Chemistry Reference Style

%%=========================================================================================%%
%% the documentclass is set to pdflatex as default. You can delete it if not appropriate.  %%
%%=========================================================================================%%

%%\documentclass[sn-basic]{sn-jnl}% Basic Springer Nature Reference Style/Chemistry Reference Style

%%Note: the following reference styles support Namedate and Numbered referencing. By default the style follows the most common style. To switch between the options you can add or remove “Numbered” in the optional parenthesis. 
%%The option is available for: sn-basic.bst, sn-chicago.bst%  
 
%%\documentclass[pdflatex,sn-nature]{sn-jnl}% Style for submissions to Nature Portfolio journals
%%\documentclass[pdflatex,sn-basic]{sn-jnl}% Basic Springer Nature Reference Style/Chemistry Reference Style
\documentclass[pdflatex,sn-mathphys-num]{sn-jnl}% Math and Physical Sciences Numbered Reference Style
%%\documentclass[pdflatex,sn-mathphys-ay]{sn-jnl}% Math and Physical Sciences Author Year Reference Style
%%\documentclass[pdflatex,sn-aps]{sn-jnl}% American Physical Society (APS) Reference Style
%%\documentclass[pdflatex,sn-vancouver-num]{sn-jnl}% Vancouver Numbered Reference Style
%%\documentclass[pdflatex,sn-vancouver-ay]{sn-jnl}% Vancouver Author Year Reference Style
%%\documentclass[pdflatex,sn-apa]{sn-jnl}% APA Reference Style
%%\documentclass[pdflatex,sn-chicago]{sn-jnl}% Chicago-based Humanities Reference Style

%%%% Standard Packages
%%<additional latex packages if required can be included here>

\usepackage{graphicx}%
\usepackage{multirow}%
\usepackage{amsmath,amssymb,amsfonts}%
\usepackage{amsthm}%
\usepackage{mathrsfs}%
\usepackage[title]{appendix}%
\usepackage{xcolor}%
\usepackage{textcomp}%
\usepackage{manyfoot}%
\usepackage{booktabs}%
\usepackage{algorithm}%
\usepackage{algorithmicx}%
\usepackage{algpseudocode}%
\usepackage{listings}%
\usepackage{subfiles}
\usepackage{threeparttable}
%%%%

%%%%%=============================================================================%%%%
%%%%  Remarks: This template is provided to aid authors with the preparation
%%%%  of original research articles intended for submission to journals published 
%%%%  by Springer Nature. The guidance has been prepared in partnership with 
%%%%  production teams to conform to Springer Nature technical requirements. 
%%%%  Editorial and presentation requirements differ among journal portfolios and 
%%%%  research disciplines. You may find sections in this template are irrelevant 
%%%%  to your work and are empowered to omit any such section if allowed by the 
%%%%  journal you intend to submit to. The submission guidelines and policies 
%%%%  of the journal take precedence. A detailed User Manual is available in the 
%%%%  template package for technical guidance.
%%%%%=============================================================================%%%%

%% as per the requirement new theorem styles can be included as shown below
\theoremstyle{thmstyleone}%
\newtheorem{theorem}{Theorem}%  meant for continuous numbers
%%\newtheorem{theorem}{Theorem}[section]% meant for sectionwise numbers
%% optional argument [theorem] produces theorem numbering sequence instead of independent numbers for Proposition
\newtheorem{proposition}[theorem]{Proposition}% 
%%\newtheorem{proposition}{Proposition}% to get separate numbers for theorem and proposition etc.

\theoremstyle{thmstyletwo}%
\newtheorem{example}{Example}%
\newtheorem{remark}{Remark}%

\theoremstyle{thmstylethree}%
\newtheorem{definition}{Definition}%

\raggedbottom
% Track-changes highlighting disabled for camera-ready: macro is now a no-op.
\newcommand{\revisiontwo}[1]{#1}

\begin{document}

\title[Real-Time Hallucination Correction in Vision-Language Models Using Dynamic Knowledge Graph Verification]{Real-Time Hallucination Correction in Vision-Language Models Using Dynamic Knowledge Graph Verification}

% ===== Authors =====

\author*[1]{\fnm{Quang-Vinh} \sur{Dang}}\email{vinh.dq4@buv.edu.vn}

\author[2]{\fnm{Ngoc-Son-An} \sur{Nguyen}}\email{annns25871@pgr.iuh.edu.vn}

\author[3]{\fnm{Thi-Bich-Diem} \sur{Vo}}\email{misdiem@gmail.com}

\author[5]{\fnm{Minh Ngoc} \sur{Dinh}}\email{minh.dinh@maeducation.com}

\author*[4]{\fnm{Anh Dat} \sur{Le}}\email{datla@uef.edu.vn}

\author[1]{\fnm{Hoang Viet} \sur{Vu}}\email{30066555@st.buv.edu.vn}

\author[1]{\fnm{Phuong Lan} \sur{Nguyen}}\email{30066017@st.buv.edu.vn}

% ===== Affiliations (based on “Added institutions” / primary affiliations) =====

\affil*[1]{\orgname{British University Vietnam}, \orgaddress{\city{Hung Yen}, \country{Viet Nam}}}
\affil[2]{\orgname{Industrial University of Ho Chi Minh City}, \orgaddress{\city{Ho Chi Minh City}, \country{Viet Nam}}}
\affil[3]{\orgname{GiaoHangNhanh}, \orgaddress{\city{Ho Chi Minh City}, \country{Viet Nam}}}
\affil[4]{\orgname{University of Economics and Finance}, \orgaddress{\city{Ho Chi Minh City}, \country{Vietnam}}}
\affil[5]{\orgname{School of Data Science and Information Technology, Millennia University Vietnam}, \orgaddress{\city{Ho Chi Minh City}, \country{Vietnam}}}


%%==================================%%
%% Sample for unstructured abstract %%
%%==================================%%

\abstract{Hallucinations in Vision-Language Models (VLMs) --- where models generate plausible but visually ungrounded or factually incorrect content --- remain a significant barrier to reliable multimodal AI deployment. Existing mitigation strategies either require expensive model retraining or rely on static knowledge bases that cannot adapt to unseen visual contexts. We present CMVKG-Guard, a training-free middleware framework that intercepts and corrects hallucinations token-by-token during autoregressive generation, without modifying the underlying VLM. The framework dynamically constructs a per-inference Hierarchical Multimodal Knowledge Graph (DH-MMKG) from visual scene entities and enriches it via live queries to Wikidata and ConceptNet. Each candidate token is scored by a three-layer Cross-Modal Verification Engine (CMVE) combining visual-textual alignment, knowledge graph grounding, and multi-hop reasoning consistency; tokens falling below an adaptive, step-aware threshold trigger a Constrained Look-Ahead Beam Search that selects a grounded replacement and teacher-forces it into the decoding sequence. Experiments on six standard benchmarks --- including POPE, CHAIR, and MMHalBench --- against eight baselines show that CMVKG-Guard reduces object hallucination rates by 84\% on CHAIR, achieves 91.5\% F1 on POPE (outperforming the strongest training-free baseline by 5.8 percentage points), and incurs a median end-to-end latency of 12ms per token on a single A100 GPU. These results suggest that dynamic, inference-time knowledge grounding is a promising direction for hallucination mitigation in open-domain VLM settings.}

\keywords{Hallucination Mitigation, Multimodal Knowledge Graphs, Cross-Modal Verification, Real-Time Correction, Adaptive Confidence Calibration}

%%\pacs[JEL Classification]{D8, H51}

%%\pacs[MSC Classification]{35A01, 65L10, 65L12, 65L20, 65L70}

\maketitle

\subfile{1intro}
\subfile{2relate}
\subfile{3method}
\subfile{4exp}
\subfile{5conclusion}
\section*{Declarations}

\textbf{Ethical approval:} Not applicable. This study is a computer science study and did not involve human participants, human data, animals, or any procedures requiring ethical approval.

\textbf{Consent to participate:} Not applicable. This study did not involve human participants.

\textbf{Consent to publish:} Not applicable. This study did not involve human participants or identifiable personal data.

\textbf{Clinical trial registration:} Not applicable. This study is not a clinical trial.

\textbf{Funding:} This research received no external funding.

\textbf{Competing interests:} The authors declare that they have no competing interests. This manuscript is not a dual publication: it has not been published previously and is not under consideration for publication elsewhere.

\textbf{Data availability:} All datasets analysed during the current study are publicly available benchmark datasets obtained from their original public repositories, and are cited in this article: POPE~\cite{li2023pope}, CHAIR (built on the MS-COCO dataset)~\cite{rohrbach2018object}, MMHalBench~\cite{sun2024mmhal}, MedHallBench~\cite{zuo2024medhallbench}, InfoSeek~\cite{chen2024extracting}, and KVQA~\cite{shah2019kvqa}. No new datasets were generated during the current study.

\backmatter

% \bmhead{Supplementary information}

% If your article has accompanying supplementary file/s please state so here. 

% Authors reporting data from electrophoretic gels and blots should supply the full unprocessed scans for key as part of their Supplementary information. This may be requested by the editorial team/s if it is missing.

% Please refer to Journal-level guidance for any specific requirements.

% \bmhead{Acknowledgements}

% Acknowledgements are not compulsory. Where included they should be brief. Grant or contribution numbers may be acknowledged.

% Please refer to Journal-level guidance for any specific requirements.

% \section*{Declarations}

% Some journals require declarations to be submitted in a standardised format. Please check the Instructions for Authors of the journal to which you are submitting to see if you need to complete this section. If yes, your manuscript must contain the following sections under the heading `Declarations':

% \begin{itemize}
% \item Funding
% \item Conflict of interest/Competing interests (check journal-specific guidelines for which heading to use)
% \item Ethics approval and consent to participate
% \item Consent for publication
% \item Data availability 
% \item Materials availability
% \item Code availability 
% \item Author contribution
% \end{itemize}

% \noindent
% If any of the sections are not relevant to your manuscript, please include the heading and write `Not applicable' for that section. 

% %%===================================================%%
% %% For presentation purpose, we have included        %%
% %% \bigskip command. Please ignore this.             %%
% %%===================================================%%
% \bigskip
% \begin{flushleft}%
% Editorial Policies for:

% \bigskip\noindent
% Springer journals and proceedings: \url{https://www.springer.com/gp/editorial-policies}

% \bigskip\noindent
% Nature Portfolio journals: \url{https://www.nature.com/nature-research/editorial-policies}

% \bigskip\noindent
% \textit{Scientific Reports}: \url{https://www.nature.com/srep/journal-policies/editorial-policies}

% \bigskip\noindent
% BMC journals: \url{https://www.biomedcentral.com/getpublished/editorial-policies}
% \end{flushleft}

% \begin{appendices}

% \section{Section title of first appendix}\label{secA1}

% An appendix contains supplementary information that is not an essential part of the text itself but which may be helpful in providing a more comprehensive understanding of the research problem or it is information that is too cumbersome to be included in the body of the paper.

% %%=============================================%%
% %% For submissions to Nature Portfolio Journals %%
% %% please use the heading ``Extended Data''.   %%
% %%=============================================%%

% %%=============================================================%%
% %% Sample for another appendix section			       %%
% %%=============================================================%%

% %% \section{Example of another appendix section}\label{secA2}%
% %% Appendices may be used for helpful, supporting or essential material that would otherwise 
% %% clutter, break up or be distracting to the text. Appendices can consist of sections, figures, 
% %% tables and equations etc.

% \end{appendices}

%%===========================================================================================%%
%% If you are submitting to one of the Nature Portfolio journals, using the eJP submission   %%
%% system, please include the references within the manuscript file itself. You may do this  %%
%% by copying the reference list from your .bbl file, paste it into the main manuscript .tex %%
%% file, and delete the associated \verb+\bibliography+ commands.                            %%
%%===========================================================================================%%

\bibliography{sn-bibliography}% common bib file
%% if required, the content of .bbl file can be included here once bbl is generated
%%\input sn-article.bbl

\end{document}
